% System Design and Technical Specification

% ----- 5.1 Technology Stack -----
\subsection{Technology Stack}
\label{subsec:stack}

The contract is Solidity 0.8.20 against OpenZeppelin Contracts v4, deployed to
Polygon PoS for its transaction costs and EVM compatibility; the trade-offs among the
scaling approaches an EVM deployment can choose between are surveyed by Gangwal
et al.~\cite{gangwal2023layer2}. Metadata is
content-addressed on IPFS through a pinning service; the issuance service is a
Node.js backend using ethers.js and the pinning provider's SDK. The client an
animator sees is a progressive web application for tracking cuts and deadlines.

Two notes on the stack rather than the design. OpenZeppelin v5 has since replaced
the \texttt{\_beforeTokenTransfer} hook this contract overrides with a unified
\texttt{\_update} hook and removed the \texttt{Counters} utility, so a production
deployment should port to it; the soulbinding logic moves across unchanged. And
the walletless experience the design assumes rests on account abstraction~\cite{buterin2021erc4337} that the prototype does not implement---it uses ordinary externally
owned accounts, which is why the friction claim of \cref{sec:evaluation} is a
property of the design rather than of what we built.

% ----- 5.2 Smart Contract -----
\subsection{Smart Contract Design}
\label{subsec:contract}

The \texttt{AnimatorSBT} contract is designed to extend OpenZeppelin's ERC-721 implementation
with two key modifications: \textbf{soulbinding} and \textbf{role-based access control}.

\subsubsection{Soulbinding Mechanism}

ERC-721 routes every token movement---mint, transfer and burn---through a single
internal hook. The contract overrides it and reverts whenever both the source and the
destination are non-zero addresses, which admits minting and burning and refuses
everything between them. \texttt{approve()} and \texttt{setApprovalForAll()} revert
unconditionally, so no approval-based path around the hook remains, and the token is
non-transferable through any ERC-721--compliant interface.

ERC-5192~\cite{daubenschuetz2022erc5192} solves an adjacent problem: it is a
signalling interface, letting tooling discover that a token is non-transferable,
where the hook above is enforcement. The two are complementary and the v2 contract
provides both---\texttt{locked()} returning true for every token, the
\texttt{Locked} event at mint, and the interface id advertised through
\texttt{supportsInterface}---while the v1 contract on Amoy has only the
enforcement.

\subsubsection{Roles and External Interface}

Access control uses two roles. \texttt{DEFAULT\_ADMIN\_ROLE} grants and revokes the
issuer role and is held by the system administrator; \texttt{ISSUER\_ROLE} mints and
revokes certificates and is held by the issuance service's operational wallet. The
separation means a compromised issuance key can issue and invalidate certificates but
cannot change who holds either role, which is what leaves the remedy of
\cref{subsec:trust} available.

The external interface is small. \texttt{mint} issues one certificate against a
recipient address and a metadata URI and \texttt{mintBatch} does the same for a list
of each, which is how certificates are issued at the end of a production.
\texttt{revoke} marks a certificate invalid without burning it, so that invalidating a
record does not erase the fact that it existed. \texttt{tokensOfOwner},
\texttt{isRevoked} and \texttt{issuedAt} are views, letting a verifier resolve an
address to the certificates it holds and check each one's status and issuance time.

\subsubsection{Gas Cost Analysis}

\Cref{tab:gas} reports the gas consumption of key operations measured in two
environments: the local Hardhat EVM (solc 0.8.20, optimizer enabled, 200 runs) and
the live Polygon Amoy testnet. \texttt{deploy} and \texttt{revoke} agree exactly.
The \texttt{mint()} figures do not, and the reason is not the environment but the
argument: \textbf{mint gas is dominated by the length of the \texttt{tokenURI}
string}, which the contract writes to storage in 32-byte words. Holding everything
else fixed and varying only the URI, a 25-character placeholder costs
127{,}579~gas, a CIDv0 pin (53 characters) 172{,}692, and a CIDv1 pin
(66 characters) 195{,}423; beyond that the curve flattens until the next word
boundary. The on-chain \texttt{mint()} carried a real CID and the earlier local
measurement a short placeholder, which accounts for the entire gap. Re-measured
locally with the CIDv1 URI the issuance service actually produces, a mint to a new
holder costs 195{,}423~gas, within 3\% of the 189{,}792 observed on-chain.

This matters beyond reconciling two columns. The one cost lever an implementer
controls is metadata addressing: a deployment that stores a CIDv1 pointer per
certificate pays \textbf{53\% more gas than one that stores a short
identifier} and resolves the pointer off-chain. It also means any reported
per-certificate cost is uninterpretable unless the URI scheme is stated, and we
state ours.

The same effect makes the on-chain \texttt{mintBatch(10)} figure an understatement:
that transaction was sent with 19-character placeholder URIs. Repeated with real
CIDv1 pins the batch costs 1{,}704{,}591~gas, or 170{,}459 per certificate. Batching
is therefore worth about 13\%, not the near-halving the placeholder measurement
suggests. Every figure in this subsection is regenerated by
\texttt{scripts/measure\_gas\_local.js}, which records the compiler settings and
the URI used alongside the measurements.

\begin{table*}[ht]
\centering
\caption{Gas consumption and estimated cost on Polygon (30~gwei, POL = \$0.22,
March 2026). ``Local'' is the Hardhat EVM measured with the CIDv1 metadata URI the
issuance service produces (solc 0.8.20, optimizer 200 runs); ``Amoy'' is the value
observed on-chain (transaction hashes in \cref{app:testnet}). The USD column is
computed from the local figures, which are reproducible from the repository. These
are \emph{projections}: they apply a fixed gas price and POL/USD rate to measured
gas used, and exclude any mainnet priority or L1 fees
(see \cref{subsec:eval_limitations}).}
\label{tab:gas}
\footnotesize
\setlength{\tabcolsep}{3pt}
\begin{tabular}{@{}lrrrr@{}}
\toprule
\textbf{Operation} & \textbf{Gas (local)} & \textbf{Gas (Amoy)} & \textbf{POL} & \textbf{USD} \\
\midrule
Deploy                 & 1{,}824{,}117 & 1{,}824{,}117 & 0.0547 & \$0.0120 \\
\texttt{mint()}        &   195{,}423 & 189{,}792$^{\dagger}$ & 0.0059 & \$0.0013 \\
\texttt{mintBatch(10)} & 1{,}704{,}591 & 1{,}025{,}449$^{\ddagger}$ & 0.0511 & \$0.0113 \\
\texttt{coSign()}      &    96{,}997 & --- & 0.0029 & \$0.0006 \\
\texttt{revoke()}      &    50{,}027 &  50{,}027 & 0.0015 & \$0.0003 \\
\bottomrule
\end{tabular}

\vspace{2pt}
\raggedright\scriptsize
$^{\dagger}$~Sent with a real CID; agrees with the local figure to within 3\%.
$^{\ddagger}$~Sent with 19-character placeholder URIs and therefore an
understatement; see the text.
\end{table*}

At these costs a self-attested contribution certificate costs approximately
\textbf{\$0.0013} on-chain, and a studio-attested one---a mint followed by a
co-signature---approximately \$0.0019. Batching ten certificates brings the
per-certificate figure down to \$0.0011. Deployment is a one-time \$0.012.
On-chain cost is thus negligible at any plausible scale, which is what makes the
off-chain storage dependency, rather than the chain, the binding constraint
(\cref{subsec:cost_sim}). \Cref{subsec:cost_sim} varies the gas price and the
token price and finds the conclusion unchanged.

% ----- 5.3 Prototype Implementation -----
\subsection{Prototype Implementation}
\label{subsec:prototype}

To demonstrate technical feasibility and on-chain correctness, we implemented a
prototype consisting of three components:
the smart contract deployed on a public testnet, an Issuance Service, and a Verification Portal.

\subsubsection{Testnet Deployment}

The \texttt{AnimatorSBT} contract was deployed on the Polygon Amoy testnet
(Chain ID: 80002) at address
\hexid{0x2cC31B9EFBD5dDfB5456EE4c76A0fcBCA5EEE117}
(block 35{,}058{,}633).\footnote{The deployment transaction and all subsequent
interactions are publicly verifiable at
\url{https://amoy.polygonscan.com/address/0x2cC31B9EFBD5dDfB5456EE4c76A0fcBCA5EEE117}.}
Seventy unit tests pass: 35 against the deployed v1 contract, 17 against the v2
co-signing reference contract, 7 against the canonical serialisation, 10 against the
disclosure commitment of \cref{subsec:privacy}, and one local end-to-end issuance run.
The suite exercises contract mechanics, which are the most standard part of the
system, and one thing that is not: the evidence hash is computed over a
recursively key-sorted serialisation, and seven of the tests confirm that it is
independent of key insertion order, stable across runs, and sensitive to any value
change. That closes a determinism gap in the earlier design, under which the same
claim serialised differently produced different hashes. \Cref{subsec:eval_limitations}
states what the suite does not reach.

An end-to-end demonstration of the Tanaka scenario
(\cref{subsec:flow}) was executed on the testnet,
successfully minting an SBT with IPFS-hosted metadata
and verifying all on-chain state (balance, token URI, revocation status).

\subsubsection{Issuance Service}

A Node.js-based Issuance Service was implemented using ethers.js~v6
and the Pinata SDK.
The service performs the complete issuance workflow:
(1)~constructing metadata JSON conforming to the schema in \cref{tab:metadata},
(2)~computing the \texttt{evidence\_hash} (Keccak-256 of the serialized attributes),
(3)~uploading the metadata to IPFS via Pinata,
and (4)~calling \texttt{AnimatorSBT.mint()} with the resulting IPFS~URI.

An end-to-end demonstration script recreates the Tanaka scenario
described in \cref{subsec:flow}:
a key animator completes 25~cuts on Episode~7, and the system
issues an SBT recording the contribution with a verifiable evidence hash.

\subsubsection{Verification Portal}

A lightweight Verification Portal was implemented as a static HTML page
using ethers.js (loaded from CDN), requiring no server-side infrastructure.
The portal accepts a wallet address, queries the contract's \texttt{tokensOfOwner()}
function, retrieves each SBT's metadata from IPFS,
and displays the contribution details including project, role, task count, and dates.
The portal also supports evidence hash verification by recomputing
the Keccak-256 hash from the IPFS metadata and comparing it
to the stored \texttt{evidence\_hash} field.

\subsubsection{Co-signing Reference Contract and Canonical Hashing}
\label{subsec:canonical}
\label{subsubsec:v2}

The authenticity layer of \cref{subsec:cosigning} exists as a reference
implementation rather than a sketch. It is validated in the local EVM, where
co-signing, canonical hashing, minting and soulbinding compose correctly; it is not
deployed on chain and does not replace the v1 contract on Amoy. \texttt{AnimatorSBTV2} adds a studio
attester role and an idempotent \texttt{coSign} taking the attester from
\texttt{msg.sender}, together with the views that report who has co-signed and at
what attestation level, and ERC-5192 \texttt{locked()} signalling. A co-signature
costs 96{,}997~gas and the contract deploys for 2{,}058{,}959.

The canonical serialisation module is the component that makes the evidence hash
mean anything. \Cref{fig:canonical} shows what it buys, with the hashes the
implementation produces. It is a dependency-light, recursively key-sorted JSON serialiser---a
subset of JCS (RFC~8785)---over which \texttt{evidence\_hash} is computed, so that
two parties assembling the same attributes in different key order arrive at the
same hash and a verifier can tell a reordering from a change. Without it the hash
fixes one serialisation of a claim rather than the claim, which was a gap in the
earlier design. The issuance service extends the v1 flow with this hash and an
optional co-signing step using a separate attester wallet, behind a
wallet-provisioning seam that is currently an unconfigured skeleton.

\begin{figure*}[t]
  \centering
  \includegraphics[width=0.92\textwidth]{figures/canonical_hash}
  \caption{The evidence hash fixes the claim rather than one spelling of it. The
  same attributes serialised with keys in different orders reduce to the same
  digest; changing a single value does not. Hashes are the values produced by the
  implementation in the accompanying repository.}
  \label{fig:canonical}
\end{figure*}

% ----- 5.4 Integration with Production Management Client -----
\subsection{Integration and the Attestation Anchor}
\label{subsec:integration}

Where a certificate is anchored is a design choice, and it is the one that decides
what a co-signature is worth. Anchoring on a self-reported work log means the
studio, asked to countersign, must review work and reach a judgment it had not
otherwise recorded. Anchoring on an accepted order means it has already reached
that judgment: taking delivery of a set of cuts and paying for them is a
determination about who did the work, made for commercial reasons and recorded in
a system the studio keeps for its own purposes. A certificate issued per accepted
order, keyed on the order identifier, is idempotent by construction, and the
co-signature of \cref{subsec:cosigning} then records a determination rather than
soliciting one. We take the anchor, and not the token, to be the part of the design that carries
the trust, and we argue it here on those grounds rather than demonstrating it.

The pattern is not new, only newly applied. Triple-entry accounting adds a third,
cryptographically verifiable entry witnessing a transaction that two parties have
already recorded in their own books, and Sarwar et al.\ survey its blockchain
implementations for business-to-business
settlement~\cite{sarwar2023tea}. A certificate anchored to an accepted order is
that third entry, with the transaction being the commissioning of work rather than
the sale of goods, and the party the entry protects being the one whose name the
books record but the credits do not. Reading the design this way also names what
it inherits: a third entry is only as good as the two it witnesses, so a studio
that mis-records who did the work produces a faithful certificate of a false
premise.

The record the anchor needs is one the commissioning party is already obliged to keep.
Article~7 of the statute governing payment to small and medium-sized contractors
requires a commissioning party to create and retain for two years a per-transaction
document recording, among sixteen enumerated items, the content of the work ordered,
the date delivery was due and the date it was received, the completion date and result
of any inspection together with the treatment of anything that failed it, and the fee
with its due date and the date and means of
payment~\cite{meti2026animeguideline}.\footnote{Act No.~120 of 1956, renamed with
effect from January 2026; the enumeration is Figure~33 of METI's guideline for this
industry, at p.~48. The duty attaches to whoever places the order: the guideline states
that a commission may fall under the Act whatever the commissioning party's own
position in the subcontracting chain.} That is the accepted order in the state the
anchor requires, held by the party whose determination a co-signature would record. The
design does not ask the industry to begin keeping a record. It asks that a record the
law already compels be made durable in the worker's hands.

The same enumeration shows why the gap it addresses exists. Nothing in the list requires
the individual creator's name---the counterparty may be recorded by number or code---and
neither that provision nor the Cabinet Office and Fair Trade Commission's 2026
guidelines on transactions at anime production sites contain any provision on crediting
or on moral rights~\cite{meti2026animeguideline,jftc2026shishin}. Regulation of this
industry has become detailed about the terms on which work is commissioned while saying
nothing about who is recorded as having done it. The data exists and is retained by
legal compulsion; what no rule produces is an attribution the worker can carry.

Two consequences follow that a self-reported log does not raise. An accepted order
can be adjusted afterwards---in Japanese practice through a correcting credit
note---including after a certificate referring to it has been issued, so an
implementation cannot rest on the token's immutability and must detect that the
underlying record changed. And the anchor presupposes a system the studio itself
operates, which a tool holding an animator's data on their own device cannot be;
the choice of anchor therefore constrains where such a system can live.

We are precise about what is and is not evaluable here. The issuance service in
the accompanying repository builds metadata from supplied attributes and is not
keyed on an order, so the anchor described above is a design argument that the
open artifact does not implement; what the repository demonstrates is the
contract, the co-signing mechanism, the canonical hash and the costs. An
order-keyed issuance path does exist behind a feature flag in a commercial
production-management system, which establishes that the integration is
constructible and nothing more. That system is not open source, no certificates
have been issued from production data, and we exclude it from the evidence this
paper rests on.
